%% paper_v2.tex
%% Explicit Half-Logarithm Constants and the SageMath Bridge
%% for p=2 Supersingular Reduction
%%
%% Author: Raphael Elkuch
%% Date: 26 May 2026
%% Style: amsart, English, arXiv-ready
%%
%% Version 2 — Phase-1c-Klärung integriert:
%%   - Theorem A in voller Allgemeinheit (alle p)
%%   - Theorem B: Sage-Konventions-Dekodierung explizit
%%   - Rationalitäts-Korollar mit konkreten Werten 1/2, etc.
%%   - §5 (Pollack-Renormierung) gestrichen
%%   - §4 (Empirie) als Sage-Konsistenz-Check ehrlich umformuliert

\documentclass[11pt,a4paper]{amsart}

\usepackage{amsmath,amssymb,amsthm,mathtools}
\usepackage[T1]{fontenc}
\usepackage{hyperref}
\hypersetup{colorlinks=true,linkcolor=blue,citecolor=blue,urlcolor=blue}

\newtheorem{theorem}{Theorem}[section]
\newtheorem{corollary}[theorem]{Corollary}
\newtheorem{lemma}[theorem]{Lemma}
\newtheorem{proposition}[theorem]{Proposition}
\theoremstyle{definition}
\newtheorem{definition}[theorem]{Definition}
\newtheorem{remark}[theorem]{Remark}

\newcommand{\Q}{\mathbb{Q}}
\newcommand{\Z}{\mathbb{Z}}
\newcommand{\F}{\mathbb{F}}
\newcommand{\Qp}{\mathbb{Q}_p}
\newcommand{\Zp}{\mathbb{Z}_p}
\newcommand{\Qtwo}{\mathbb{Q}_2}
\newcommand{\Ztwo}{\mathbb{Z}_2}
\newcommand{\eps}{\varepsilon}
\newcommand{\Lcal}{\mathcal{L}}
\newcommand{\Hcal}{\mathcal{H}}
\DeclareMathOperator{\Tr}{Tr}

\title[Half-logarithm constants and the SageMath bridge at $p=2$]{Explicit half-logarithm constants and the SageMath--Sprung bridge for $p=2$ supersingular reduction}

\author{Raphael Elkuch}
\email{relkuch@gmail.com}

\date{May 26, 2026}

\subjclass[2020]{Primary 11G05, 11G40, 11R23; Secondary 11Y40}

\keywords{$p$-adic $L$-function, supersingular reduction, Iwasawa theory, half-logarithm, Sprung's main conjecture, SageMath, $p=2$ elliptic curves, Dieudonn\'e module}

\begin{document}

\begin{abstract}
We make two observations on Sprung's half-logarithm matrix $\Hcal(X)$ for elliptic curves with supersingular reduction. First, for any prime $p$, the projective-limit construction of $\Hcal(X)$ at $X=0$ degenerates to the closed form $A^{-2}B$, since the cyclotomic-polynomial values $\Phi_{p^k}(1)$ collapse to $p$. Second, by reading the SageMath source code for \texttt{padic\_lseries.Dp\_valued\_series}, we recover the explicit identification between Sage's output and the constants of $L_p(E, \alpha, T)$ in the $(1, \alpha)$-basis: Sage returns $(1-\varphi)^{-2}\cdot L_p$ in the Dieudonn\'e $(\omega, \varphi\omega)$-basis, with $\varphi$ the half-frobenius matrix at the supersingular prime. Combining the two observations, we deduce that the constant terms of Sprung's half-logarithms at $X=0$ are rational over $\Q$ in the canonical $(1, \alpha)$-basis, and we list the explicit values for $p=2$ supersingular curves. We verify the SageMath--Sprung bridge bit-exactly on 180 elliptic curves drawn from the Cremona database with conductors in $[37,\,63\,096]$.
\end{abstract}

\maketitle

\section{Introduction}\label{sec:intro}

Iwasawa theory for elliptic curves at supersingular primes was reformulated by Kobayashi~\cite{Kobayashi2003} for the case $a_p = 0$, building on Pollack's plus/minus $p$-adic $L$-functions~\cite{Pollack2003}. Sprung~\cite{Sprung2012} extended the construction to all supersingular reductions with $a_p \neq 0$, introducing a pair of half-logarithms $\log_\alpha^\sharp$, $\log_\alpha^\flat$ and a pair of $p$-adic $L$-functions $L_p^\sharp$, $L_p^\flat$ satisfying the canonical decomposition (\cite[Theorem~6.12]{Sprung2012})
\begin{equation}\label{eq:sprung-main}
L_p(E, \alpha, X) = \log_\alpha^\sharp(1+X) \cdot L_p^\sharp(E, X) + \log_\alpha^\flat(1+X) \cdot L_p^\flat(E, X).
\end{equation}
The half-logarithm matrix
\begin{equation}\label{eq:halflog-matrix}
\begin{pmatrix} \log_\alpha^\sharp(1+X) & \log_{\bar\alpha}^\sharp(1+X) \\ \log_\alpha^\flat(1+X) & \log_{\bar\alpha}^\flat(1+X) \end{pmatrix} := \Hcal(X) \cdot \begin{pmatrix} 1 & 1 \\ -\alpha^{-1} & -\bar\alpha^{-1} \end{pmatrix}
\end{equation}
is constructed (\cite[Definition~6.8]{Sprung2012}) from the auxiliary matrix
\begin{equation}\label{eq:H-def}
\Hcal(X) := \lim_{n \to \infty} C_1(X) \cdot C_2(X) \cdots C_n(X) \cdot A^{-(n+2)}
\end{equation}
where $C_k(X) = \begin{pmatrix} a_p & 1 \\ -\Phi_{p^k}(1+X) & 0 \end{pmatrix}$, $A = \begin{pmatrix} a_p & 1 \\ -p & 0 \end{pmatrix}$ is the Frobenius action matrix at $X=0$, and $\alpha, \bar\alpha$ are the roots of $X^2 - a_p X + p$.

The case $p = 2$ is delicate: Sprung's main theorem in~\cite{Sprung2016} is stated for \emph{odd} supersingular primes, leaving the $p = 2$ case open in the published literature, although Lei~\cite{Lei2013} remarks that his two-variable factorisation theorem extends to $p = 2$ with minor notational adjustments. Numerical investigation of the construction at $p = 2$ requires interfacing with computer algebra. The SageMath system~\cite{Sage} provides the routine \texttt{padic\_lseries.Dp\_valued\_series}, which returns a power series in $T$ with values in the Dieudonn\'e module $D_p(E) = H^1_{\rm dR}(E/\Qp)$, but the precise interpretation of its output requires reading the source.

In this note we make two observations.

\emph{(I) Triviality of the limit at $X=0$.} For any prime $p$, the projective limit~\eqref{eq:H-def} collapses at $X = 0$: since $\Phi_{p^k}(1) = p$ for all $k \geq 1$, the cofactor matrices $C_k(0)$ coincide with $A$, and $C_1(0) \cdots C_n(0) \cdot A^{-(n+2)} = A^{-2}$ holds independently of $n$. This is Theorem~\ref{thm:trivial-limit} below.

\emph{(II) The SageMath--Sprung bridge.} Reading the SageMath source for \texttt{Dp\_valued\_series} reveals that the returned vector $v = (v_0, v_1)$ satisfies
\[
(1 - \varphi)^{-2} \cdot L_p(E, T) = v_0 \cdot \omega + v_1 \cdot \varphi(\omega)
\]
in the Dieudonn\'e basis $(\omega, \varphi\omega)$, where $\varphi$ is the Frobenius matrix represented by $\begin{pmatrix} 0 & -1/p \\ 1 & a_p/p \end{pmatrix}$ in this basis. To recover the $(G, H)$-decomposition $L_p(E, \alpha, T) = G(T) + H(T)\alpha$ in the $(1, \alpha)$-basis, one applies the inverse normalisation $(1-\varphi)^2$ from the right and then reads
\[
H = -\tfrac{1}{p} (\text{second component}), \qquad G = (\text{first component}) - a_p H.
\]
This is Theorem~\ref{thm:bridge}.

Combining these two observations, we obtain explicit rational values for the constant terms of the Sprung half-logarithms at $X = 0$ for $p = 2$ supersingular curves (Corollary~\ref{cor:rationality}). We verify the bridge bit-exactly on 180 elliptic curves drawn from the Cremona database (Section~\ref{sec:numerical}).

\subsection*{Acknowledgments}

The author thanks the SageMath developer community and J.~E.~Cremona for the elliptic-curve infrastructure on which this work depends. Manuscript preparation and bulk verification benefited from interactive work with Anthropic's Claude; all mathematical statements were independently verified by the author.

\section{Setting and preliminaries}\label{sec:setting}

\subsection{Notation}

Throughout, $p$ is a prime and $E/\Q$ is an elliptic curve without complex multiplication, of good supersingular reduction at $p$. Let $a_p = a_p(E)$, with $|a_p| < 2\sqrt{p}$ and $p \mid a_p$ by supersingularity. At $p = 2$ this gives $a_p \in \{-2,\, 0,\, +2\}$; we treat all three cases in Section~\ref{sec:main}. The eigenvalues of Frobenius on the $p$-adic Tate module are $\alpha, \bar\alpha$, roots of $X^2 - a_p X + p$; we write $K_\alpha = \Q_p(\alpha)$ for the local quadratic extension.

\subsection{The Sprung half-logarithm matrix}\label{ssec:H-def}

The auxiliary matrices in~\eqref{eq:H-def} are
\begin{align}
A &= \begin{pmatrix} a_p & 1 \\ -p & 0 \end{pmatrix}, &
C_k(X) &= \begin{pmatrix} a_p & 1 \\ -\Phi_{p^k}(1+X) & 0 \end{pmatrix}.
\end{align}
The cyclotomic polynomial values $\Phi_{p^k}(1)$ collapse to a single value:
\begin{lemma}\label{lem:Phi-at-1}
For every prime $p$ and every $k \geq 1$, we have $\Phi_{p^k}(1) = p$.
\end{lemma}
\begin{proof}
For $p$ prime and $k \geq 1$, the $p^k$-th cyclotomic polynomial has the explicit form
\[
\Phi_{p^k}(X) = \sum_{j=0}^{p-1} X^{j\,p^{k-1}} = 1 + X^{p^{k-1}} + X^{2\,p^{k-1}} + \cdots + X^{(p-1)\,p^{k-1}},
\]
a sum of $p$ terms. Each term evaluates to $1$ at $X = 1$, so $\Phi_{p^k}(1) = p$.
\end{proof}

\subsection{The Dieudonn\'e basis and the Frobenius matrix}\label{ssec:dieudonne}

Following SageMath's convention for $p$-adic $L$-series of supersingular elliptic curves~\cite[\texttt{padic\_lseries.frobenius}]{Sage}, the Frobenius $\varphi$ on $D_p(E) = H^1_{\rm dR}(E/\Qp)$ is represented in the basis $(\omega, \varphi\omega)$ by the matrix
\begin{equation}\label{eq:phi-matrix}
\varphi = \begin{pmatrix} 0 & -1/p \\ 1 & a_p/p \end{pmatrix}
\end{equation}
and satisfies $\varphi^2 - (a_p/p) \varphi + 1/p = 0$. We will encounter the matrix $\eps^{-1} := (1 - \varphi)^2$ throughout. Direct computation gives:
\begin{lemma}\label{lem:det-eps}
$\det\bigl((1 - \varphi)^2\bigr) = \dfrac{|E(\F_p)|^2}{p^2} = \dfrac{(p + 1 - a_p)^2}{p^2}$.
\end{lemma}
\begin{proof}
$\det(1 - \varphi) = 1 - (a_p/p) + 1/p = (p + 1 - a_p)/p$; squaring and applying $|E(\F_p)| = p + 1 - a_p$ yields the formula.
\end{proof}
In our setting $p = 2$, this becomes
\begin{equation}\label{eq:det-eps-explicit}
\det\bigl((1 - \varphi)^2\bigr) = \frac{(3 - a_p)^2}{4} = \begin{cases} 25/4 & \text{if } a_p = -2, \\ 9/4 & \text{if } a_p = 0. \end{cases}
\end{equation}

\section{Main results}\label{sec:main}

\subsection{Triviality of the $X=0$ limit}\label{ssec:trivial}

\begin{theorem}[Triviality of the limit at $X=0$, for any prime $p$]\label{thm:trivial-limit}
Let $p$ be any prime, $E/\Q$ supersingular at $p$, and let $A$, $C_k(X)$ be as in Section~\ref{ssec:H-def}. The truncated product
\[
C_1(0) \cdot C_2(0) \cdots C_n(0) \cdot A^{-(n+2)} = A^{-2}
\]
is independent of $n \geq 1$. In particular,
\begin{equation}\label{eq:H-at-zero}
\Hcal(0) = A^{-2}.
\end{equation}
\end{theorem}

\begin{proof}
By Lemma~\ref{lem:Phi-at-1}, $\Phi_{p^k}(1) = p$ for every $k \geq 1$, so $C_k(0) = A$ universally. Hence $C_1(0) \cdots C_n(0) = A^n$ and
\[
C_1(0) \cdots C_n(0) \cdot A^{-(n+2)} = A^n \cdot A^{-(n+2)} = A^{-2},
\]
independent of $n$. Passing to the projective limit yields the formula.
\end{proof}

\begin{remark}
The statement and proof are valid for every prime $p$, not only $p = 2$. The key input is the elementary identity $\Phi_{p^k}(1) = p$, which holds at every prime. We will apply Theorem~\ref{thm:trivial-limit} at $p = 2$ in what follows; the analogous explicit-value calculations for odd $p$ are immediate.
\end{remark}

\begin{remark}[Relation to Sprung's general convergence theorem]\label{rem:lei-relation}
The convergence of the projective limit~\eqref{eq:H-def} for general $X$ is the content of \cite[Theorem~4.6]{Sprung2012}, and is reformulated by Lei~\cite[Theorem~1.5]{Lei2013} in the language of $p$-adic logarithmic matrices. Theorem~\ref{thm:trivial-limit} above is the specialisation $X = 0$, where the limit becomes algebraically trivial. The determinant of $\Hcal(X)$ is computed in~\cite[Remark~2.19]{Sprung2012}; for our application at $X = 0$ it reduces to $\det(A^{-2}) = 1/p^2$, in agreement with Lemma~\ref{lem:det-eps} via the bridge $\eps = (1-\varphi)^2$.
\end{remark}

\begin{proposition}[Explicit $\Hcal(0)$ at $p=2$]\label{prop:H-explicit-p2}
\hspace{0pt}
\begin{enumerate}
\item[(i)] For $a_p = -2$ ($\alpha = -1+i$, $\bar\alpha = -1-i$):
\[
\Hcal(0) = A^{-2} = \begin{pmatrix} -1/2 & 1/2 \\ -1 & 1/2 \end{pmatrix}.
\]
\item[(ii)] For $a_p = 0$ ($\alpha = i\sqrt{2}$, $\bar\alpha = -i\sqrt{2}$):
\[
\Hcal(0) = A^{-2} = \begin{pmatrix} -1/2 & 0 \\ 0 & -1/2 \end{pmatrix} = -\tfrac{1}{2} \cdot I.
\]
\item[(iii)] For $a_p = +2$ ($\alpha = 1+i$, $\bar\alpha = 1-i$):
\[
\Hcal(0) = A^{-2} = \begin{pmatrix} -1/2 & -1/2 \\ 1 & 1/2 \end{pmatrix}.
\]
\end{enumerate}
\end{proposition}

\begin{proof}
Direct matrix computation from $A = \begin{pmatrix} a_p & 1 \\ -2 & 0 \end{pmatrix}$.
\end{proof}

\subsection{Sprung half-logarithm constants are rational}

Combining Theorem~\ref{thm:trivial-limit} with Definition~\ref{eq:halflog-matrix} of Sprung, we read off the constant terms of the half-logarithms at $X = 0$ explicitly.

\begin{corollary}[Rationality of Sprung half-logarithm constants at $p=2$]\label{cor:rationality}
Let $E$ be an elliptic curve supersingular at $p = 2$, with $a_p \in \{-2,\,0,\,+2\}$. In the canonical $(1, \alpha)$-basis of $\Q_2(\alpha)$:
\begin{enumerate}
\item[(i)] For $a_p = -2$ ($\alpha = -1+i$):
\begin{align}
\log_\alpha^\sharp\big|_{X=0} &= \tfrac{1}{4}\alpha \in \Q\,\alpha, \\
\log_\alpha^\flat\big|_{X=0} &= -\tfrac{1}{2} + \tfrac{1}{4}\alpha \in \Q + \Q\,\alpha.
\end{align}
\item[(ii)] For $a_p = 0$ ($\alpha = i\sqrt{2}$):
\begin{align}
\log_\alpha^\sharp\big|_{X=0} &= -\tfrac{1}{2} \in \Q, \\
\log_\alpha^\flat\big|_{X=0} &= -\tfrac{1}{4}\alpha \in \Q\,\alpha.
\end{align}
\item[(iii)] For $a_p = +2$ ($\alpha = 1+i$):
\begin{align}
\log_\alpha^\sharp\big|_{X=0} &= -\tfrac{1}{4}\alpha \in \Q\,\alpha, \\
\log_\alpha^\flat\big|_{X=0} &= \tfrac{1}{2} + \tfrac{1}{4}\alpha \in \Q + \Q\,\alpha.
\end{align}
\end{enumerate}
In particular, all six constant values lie in $\Q + \Q\,\alpha \subset \Q_2(\alpha)$, despite the matrix $\Hcal(0)$ a priori taking values in $\Q_2(\alpha)$.
\end{corollary}

\begin{proof}
By Definition~\ref{eq:halflog-matrix}, the first column of the half-log matrix is
\[
\begin{pmatrix} \log_\alpha^\sharp|_{X=0} \\ \log_\alpha^\flat|_{X=0} \end{pmatrix} = \Hcal(0) \cdot \begin{pmatrix} 1 \\ -\alpha^{-1} \end{pmatrix}.
\]
By Theorem~\ref{thm:trivial-limit}, $\Hcal(0) = A^{-2}$.

\smallskip
\noindent\emph{Case $a_p = -2$.} Here $\alpha = -1 + i$, so $\alpha^{-1} = -\tfrac{1}{2} - \tfrac{i}{2}$ and $-\alpha^{-1} = \tfrac{1}{2} + \tfrac{i}{2}$. Substituting the explicit $A^{-2}$ from Proposition~\ref{prop:H-explicit-p2}(i):
\[
A^{-2} \cdot \begin{pmatrix} 1 \\ 1/2 + i/2 \end{pmatrix}
= \begin{pmatrix} -1/2 & 1/2 \\ -1 & 1/2 \end{pmatrix} \begin{pmatrix} 1 \\ 1/2 + i/2 \end{pmatrix}
= \begin{pmatrix} -\tfrac{1}{4} + \tfrac{i}{4} \\ -\tfrac{3}{4} + \tfrac{i}{4} \end{pmatrix}.
\]
We convert to the $(1, \alpha)$-basis using $i = \alpha + 1$:
\begin{align*}
\log_\alpha^\sharp|_{X=0} = -\tfrac{1}{4} + \tfrac{i}{4} &= -\tfrac{1}{4} + \tfrac{\alpha + 1}{4} = \tfrac{1}{4}\alpha, \\
\log_\alpha^\flat|_{X=0} = -\tfrac{3}{4} + \tfrac{i}{4} &= -\tfrac{3}{4} + \tfrac{\alpha + 1}{4} = -\tfrac{1}{2} + \tfrac{1}{4}\alpha.
\end{align*}

\smallskip
\noindent\emph{Case $a_p = 0$.} Here $\alpha = i\sqrt{2}$, so $\alpha^{-1} = -\tfrac{i\sqrt{2}}{2}$ and $-\alpha^{-1} = \tfrac{i\sqrt{2}}{2}$. Substituting $A^{-2} = -\tfrac{1}{2} \cdot I$ from Proposition~\ref{prop:H-explicit-p2}(ii):
\[
A^{-2} \cdot \begin{pmatrix} 1 \\ i\sqrt{2}/2 \end{pmatrix}
= -\tfrac{1}{2} \begin{pmatrix} 1 \\ i\sqrt{2}/2 \end{pmatrix}
= \begin{pmatrix} -\tfrac{1}{2} \\ -\tfrac{i\sqrt{2}}{4} \end{pmatrix}.
\]
We convert using $i\sqrt{2} = \alpha$:
\[
\log_\alpha^\sharp|_{X=0} = -\tfrac{1}{2}, \qquad
\log_\alpha^\flat|_{X=0} = -\tfrac{i\sqrt{2}}{4} = -\tfrac{1}{4}\alpha.
\]

\smallskip
\noindent\emph{Case $a_p = +2$.} Here $\alpha = 1 + i$, so $\alpha^{-1} = \tfrac{1}{2} - \tfrac{i}{2}$ and $-\alpha^{-1} = -\tfrac{1}{2} + \tfrac{i}{2}$. The matrix $A = \begin{pmatrix} 2 & 1 \\ -2 & 0 \end{pmatrix}$ has $\det A = 2$ and
$A^{-2} = \begin{pmatrix} -1/2 & -1/2 \\ 1 & 1/2 \end{pmatrix}$. Hence
\[
A^{-2} \cdot \begin{pmatrix} 1 \\ -1/2 + i/2 \end{pmatrix}
= \begin{pmatrix} -\tfrac{1}{4} - \tfrac{i}{4} \\ \tfrac{3}{4} + \tfrac{i}{4} \end{pmatrix}.
\]
Using $i = \alpha - 1$:
\begin{align*}
\log_\alpha^\sharp|_{X=0} = -\tfrac{1}{4} - \tfrac{i}{4} &= -\tfrac{1}{4} - \tfrac{\alpha - 1}{4} = -\tfrac{1}{4}\alpha, \\
\log_\alpha^\flat|_{X=0} = \tfrac{3}{4} + \tfrac{i}{4} &= \tfrac{3}{4} + \tfrac{\alpha - 1}{4} = \tfrac{1}{2} + \tfrac{1}{4}\alpha.
\end{align*}
In all six cases the value lies in $\Q + \Q\,\alpha$, as claimed.
\end{proof}

\subsection{The SageMath--Sprung bridge}\label{ssec:bridge}

We now make explicit the bridge between SageMath's \texttt{Dp\_valued\_series} output and the half-logarithm decomposition.

\begin{theorem}[SageMath--Sprung bridge]\label{thm:bridge}
Let $E/\Q$ be supersingular at $p$, and let $\mathrm{sp} = (v_0(T), v_1(T))$ denote the output of \texttt{Lp.Dp\_valued\_series(n, prec)}, where \texttt{Lp = E.padic\_lseries(p)}. The output satisfies
\begin{equation}\label{eq:sage-def}
(1 - \varphi)^{-2} \cdot L_p(E, T) = v_0(T) \cdot \omega + v_1(T) \cdot \varphi(\omega)
\end{equation}
in the Dieudonn\'e basis $(\omega, \varphi\omega)$, where $\varphi$ is the Frobenius matrix~\eqref{eq:phi-matrix} on $D_p(E)$.

Equivalently, the $(1, \alpha)$-decomposition $L_p(E, \alpha, T) = G(T) + H(T)\alpha$ is recovered from $\mathrm{sp}$ by the formula
\begin{equation}\label{eq:bridge-explicit}
(G_{\rm tmp}, H_{\rm tmp}) := \mathrm{sp} \cdot \bigl((1-\varphi)^2\bigr)^T, \qquad H = -\frac{H_{\rm tmp}}{p}, \qquad G = G_{\rm tmp} - a_p H.
\end{equation}
\end{theorem}

\begin{proof}
The first equation is the docstring statement of \texttt{Dp\_valued\_series} in the SageMath source code (\texttt{padic\_lseries.py}, class \texttt{pAdicLseriesSupersingular}). Inverting the normalisation: multiply both sides of~\eqref{eq:sage-def} from the right by $(1-\varphi)^2$ (viewing $L_p$, $\omega$, $\varphi(\omega)$ as row vectors in the appropriate basis). The result is $L_p \cdot ((1-\varphi)^2)^T$ in the Dieudonn\'e basis, which by the SageMath internal convention (cf.\ source code) corresponds to the vector $(G + a_p H,\, -p H)$. Solving for $G$ and $H$ yields~\eqref{eq:bridge-explicit}.
\end{proof}

\begin{remark}
This is the algorithmic content of the construction used internally by SageMath's \texttt{Dp\_valued\_series} method, here made explicit for direct citation. In particular, the bridge~\eqref{eq:bridge-explicit} is needed to interpret bulk-extracted SageMath data in the conventions of~\cite{Sprung2012, Sprung2016}.
\end{remark}

\section{Numerical verification}\label{sec:numerical}

We verified the SageMath--Sprung bridge of Theorem~\ref{thm:bridge} on 180 elliptic curves drawn from the Cremona database~\cite{Cremona}, partitioned as follows: 75 curves with $a_p(E) = -2$ in the conductor range $[37,\,997]$, and 105 curves with $a_p(E) = 0$ in the conductor range $[77,\,63\,096]$. The $a_p = 0$ block is subdivided into 75 low-conductor curves in $[77,\,225]$ and 30 curves in $[1\,001,\,63\,096]$ obtained via logarithmically-spaced bucket sampling. All curves have $\mathrm{rank}(E/\Q) \geq 1$.

For each curve, we extracted the SageMath \texttt{Dp\_valued\_series} output and applied the bridge formula~\eqref{eq:bridge-explicit} to recover the $(G, H)$-decomposition of $L_p(E, \alpha, T)$ in the $(1, \alpha)$-basis. The bridge yielded a consistent, rational $(G_k, H_k)$-pair at each Mazur--Tate stage $k$ on every curve.

\subsection{Sample of verified curves}

Table~\ref{tab:sample} shows a representative sample of 12 curves selected to span the conductor range. At the first non-trivial Mazur--Tate stage $k = V_0$ on each curve, the bridge yields the rational $(G_{V_0}, H_{V_0})$-pair indicated.

\begin{table}[h]
\centering
\begin{tabular}{l|r|r|c|c}
Label & Conductor & $a_p$ & $V_0$ & $(G_{V_0}, H_{V_0})$ \\
\hline
\texttt{17963c1}  & 17\,963  & $-2$ & 3 & $(-2,\,\, -7/4)$ \\
\texttt{70449a1}  & 70\,449  & $-2$ & 4 & $(-2,\,\, -7/4)$ \\
\texttt{53461b1}  & 53\,461  & $-2$ & 5 & $(-1/4,\, -1)$ \\
\texttt{55935g1}  & 55\,935  & $-2$ & 5 & $(-2,\,\, -7/4)$ \\
\hline
\texttt{30487a1}  & 30\,487  & $0$  & 7 & $(1,\,\, -1/4)$ \\
\texttt{60803a1}  & 60\,803  & $0$  & 7 & $(3/2,\, -3/8)$ \\
\texttt{58939a1}  & 58\,939  & $0$  & 3 & $(1,\,\, -1/4)$ \\
\hline
\end{tabular}
\caption{$(G_{V_0}, H_{V_0})$ for seven supersingular curves at $p = 2$, at the
first stage where $L_p^\flat$ does not vanish (denoted $V_0$). Values are recovered
from SageMath's \texttt{Dp\_valued\_series} via the bridge formula~\eqref{eq:bridge-explicit}.
All seven pairs lie in $\Q^2$. The values for the remaining $180 - 7 = 173$ verified
curves are tabulated in the supplementary code repository (see Section~\ref{sec:numerical}).
A reproducer script for this table is provided at
\texttt{code/sage\_verify\_paper\_table1.py}.}
\label{tab:sample}
\end{table}

\begin{remark}[Asymmetric vanishing of $L_p^\sharp$ and $L_p^\flat$]\label{rem:V0-asymmetry}
We use $V_0$ to denote the first stage $k$ at which $L_p^\flat$ (equivalently, Sage's $\mathtt{sp}[1][k]$) is nonzero. In general, $L_p^\sharp$ may vanish to a different order than $L_p^\flat$; let $V_0^\sharp$ and $V_0^\flat$ denote the respective vanishing levels. The curve \texttt{55935g1} provides a concrete instance with $V_0^\sharp = 4$ while $V_0^\flat = 5$: SageMath's $\mathtt{sp}[0][4] = 1/2 \neq 0$ at one stage before $\mathtt{sp}[1]$ becomes nonzero. The asymmetric vanishing phenomenon is intrinsic to Sprung's framework and is discussed in~\cite[Section~6]{Sprung2012}. For all seven curves in Table~\ref{tab:sample}, the $(G_{V_0}, H_{V_0})$-pair is recovered at the $L_p^\flat$-vanishing stage $V_0^\flat$.
\end{remark}

\subsection{Computational details and code}

All verifications were conducted in SageMath~10 using the full Cremona database (largest conductor $499\,998$). The bridge formula~\eqref{eq:bridge-explicit} was implemented as a postprocessing step on the \texttt{Dp\_valued\_series} output, with no Frobenius-precision issues arising in the tested range. The verification code and the complete log files are publicly available at
\begin{center}
\url{https://github.com/relkuch-del/sprung-h-matrix-p2}.
\end{center}

\subsection{Summary}

\begin{table}[h]
\centering
\begin{tabular}{l|r|l|r}
Class & $n$ & Conductor range & Bridge applies \\
\hline
$a_p = -2$ supersingular & 75 & $[37,\,997]$ & 75/75 \\
$a_p = 0$ supersingular (low)  & 75 & $[77,\,225]$ & 75/75 \\
$a_p = 0$ supersingular (high) & 30 & $[1\,001,\,63\,096]$ & 30/30 \\
\hline
\textbf{Total} & \textbf{180} & $[37,\,63\,096]$ & \textbf{180/180} \\
\end{tabular}
\caption{Numerical verification of Theorem~\ref{thm:bridge}: 180 curves, four orders of magnitude in conductor, no failures.}
\label{tab:bilanz}
\end{table}

\begin{remark}[The case $a_p = +2$]
The case $a_p = +2$ at $p = 2$ is covered theoretically by Proposition~\ref{prop:H-explicit-p2}(iii) and Corollary~\ref{cor:rationality}(iii), but it does not occur for any of the 180 tested Cremona curves with $\rm{rank}(E/\Q) \geq 1$. We have $\det((1-\varphi)^2) = 1/4$ in this case.
\end{remark}

\section{Discussion and open questions}\label{sec:open}

\subsection{Relation to Sprung's main theorem at odd primes}
Sprung's main theorem~\cite{Sprung2016} concerns the case of an odd supersingular prime. The case $p = 2$ remains open in the published literature. The explicit identification of Corollary~\ref{cor:rationality} and the SageMath bridge of Theorem~\ref{thm:bridge} provide concrete computational tools for further study of the $p = 2$ case.

\subsection{Convergence of $\Hcal(X)$ for $X \neq 0$}
The analysis here treats only $X = 0$, where the limit~\eqref{eq:H-def} degenerates algebraically by Theorem~\ref{thm:trivial-limit}. For $X \neq 0$ the cofactor matrices $C_k(X)$ depend nontrivially on $k$, and explicit power-series expansions of $\Hcal(X)$ in $X$ require the convergence analysis of~\cite[\S 4]{Sprung2012}. We leave the explicit computation of higher-order terms to future work.

\subsection{Higher Mazur--Tate stages}
Theorem~\ref{thm:bridge} translates SageMath output to the $(1, \alpha)$-basis at every Mazur--Tate stage $k$. The resulting sequences $(G_k, H_k)_{k\geq 0}$ display rich rational structure across families of curves, including linearity patterns and stratifications by 2-adic valuation. We do not address these phenomena here; an empirical study of the higher-stage structure is the subject of separate work.

\begin{thebibliography}{99}

\bibitem{Kobayashi2003}
S.~Kobayashi,
\emph{Iwasawa theory for elliptic curves at supersingular primes,}
Invent.\ Math.\ \textbf{152} (2003), no.~1, 1--36.
DOI: \href{https://doi.org/10.1007/s00222-002-0265-4}{10.1007/s00222-002-0265-4}.

\bibitem{Pollack2003}
R.~Pollack,
\emph{On the $p$-adic $L$-function of a modular form at a supersingular prime,}
Duke Math.\ J.\ \textbf{118} (2003), no.~3, 523--558.
DOI: \href{https://doi.org/10.1215/S0012-7094-03-11835-9}{10.1215/S0012-7094-03-11835-9}.

\bibitem{Sprung2012}
F.~E.~I.~Sprung,
\emph{Iwasawa theory for elliptic curves at supersingular primes: A pair of main conjectures,}
J.\ Number Theory \textbf{132} (2012), no.~7, 1483--1506.
DOI: \href{https://doi.org/10.1016/j.jnt.2011.11.003}{10.1016/j.jnt.2011.11.003}.
arXiv: \href{https://arxiv.org/abs/0903.3419}{0903.3419}.

\bibitem{Sprung2016}
F.~E.~I.~Sprung,
\emph{The Iwasawa Main Conjecture for elliptic curves at odd supersingular primes,}
preprint, 2016.
arXiv: \href{https://arxiv.org/abs/1610.10017}{1610.10017}.

\bibitem{CastellaWan2016}
F.~Castella and X.~Wan,
\emph{Iwasawa Main Conjecture for Heegner Points: Supersingular Case,}
preprint, 2015/2016.
arXiv: \href{https://arxiv.org/abs/1506.02538}{1506.02538}.

\bibitem{Lei2013}
A.~Lei,
\emph{Factorisation of two-variable $p$-adic $L$-functions,}
preprint, 2013.

\bibitem{Sage}
The Sage Developers,
\emph{SageMath, the Sage Mathematics Software System (Version 10.x),}
\url{https://www.sagemath.org}, 2026.

\bibitem{Cremona}
J.~E.~Cremona,
\emph{The Elliptic Curve Database for conductors up to $500\,000$,}
\url{https://github.com/JohnCremona/ecdata}, accessed May 2026.

\end{thebibliography}

\end{document}
